\section{Secondary modification episodes}\label{sec:episodes}

Keep the reference and observer interface of \cref{sec:references}. A tape
is \emph{matched} when its current whole word differs in at most $s$ cells
from some current reference. Define a canonical carrier on the original run:
initially choose the first unmatched tape, or tape zero if all match;
subsequently retain the carrier if it is unmatched, otherwise choose the
first unmatched tape, retaining the old carrier if all match. This rule is
independent of the hash used in the analysis. Whenever any tape is unmatched,
the carrier is unmatched.

On a run with $\tau$ successful transitions, an inactive unmatched episode
of tape $h$ is a maximal interval $I=[u,v)$ of reached configuration indices
$0,\ldots,\tau$ on which $h$ is unmatched and is not the carrier. Define
\[
 D_h(I)=\{y_h(t):t\in I,\ t<\tau,\text{ the old-position write at }t
                   \text{ changes the symbol}\}.
\]
Addresses, not write occurrences, are counted. A restore still counts its
address. A successful exit transition is included, the entering transition
$u-1$ is excluded, and the halting configuration contributes no write.

\begin{lemma}[Changed cells and a frozen word]\label{lem:changed-cell-support}
\formal{HistoricalCache.lean}{ExecutionPrefix.unchanged_outside_changedCells,ExecutionPrefix.residual_card_le,ExecutionPrefix.residual_card_le_entering_add_episode,ExecutionPrefix.changedCells_card_le_duration}
In an actual successful execution prefix, let $D(u,v)$ be the distinct
old addresses changed on one tape in steps $[u,v)$. Outside $D(u,v)$
its value at $v$ equals its value at $u$. In any finite address window,
the number of cells differing from a fixed word at $v$ is at most their
number at $u$ plus $|D(u,v)|$. In particular, an initial residual of at
most $s$ cells followed by one entering step and at most $m$ distinct
changed addresses in $[u+1,v)$ leaves at most $s+m+1$ residual cells.
Also $|D(u,v)|\le v-u$.
\end{lemma}
\begin{outline}
Induct through actual successful transitions. A cell different from the
old head position is preserved; at that position it is preserved unless
the step belongs to the changed-cell set. The final residual is contained
in the union of its initial residual and this set. Split off the first
old position and use finite union cardinality. Finally the changed-cell
set is the image of a subset of $[u,v)$, giving the duration bound.
The set records actual symbol changes, so identity writes are excluded
and repeated changes or restores use only one address.
\end{outline}

\begin{theorem}[Historical reference caches]\label{thm:historical-cache}
\notformalized
Let $G$ satisfy the source/observer, halt and output-length promises of
\cref{thm:reference-carrier}, replacing its one-exception promise by
$|D_h(I)|\le m$ for every episode. Put $M=s+m+1$, retain $K_j,r,p$
from that theorem, and set
\[
 B_{\rm cache}=k(T+1)n^{k+2}(4T+1)(1+n^2)
                  [1+(T+1)(P+1)^2\ell]^k.
\]
An analysis system exact on at least half of $G$ has the same time, input
trajectory and append times/bits, with
\[
 Q_*\le2+QU B_{\rm cache}K_M^k2^{k(\ell+1)r},\qquad
 z\le Z_0+k\ell r.
\]
It uses one blank physical work tape and the same continuous ports and
radius as before. There is no restriction on episode length or count,
or on the number of simultaneous unmatched tapes.
\end{theorem}
\begin{informalproof}
Choose the same common hash event as in \cref{thm:reference-carrier}:
compare only current actual words with current radius-$s$ candidates on
original trajectories. It gives the true matching decisions on at least
half the inputs. No historical radius-$M$ search is included in the event.

Give every tape a cache consisting of a frozen base, either BLANK or
$(e,o,b,i)$, and an explicit error $\epsilon\in\operatorname{Err}_M$.
At current prefix $Y$, the latter base denotes
\[
 \mathbf1_{[0,b)}\gamma_{hi}(\,\cdot\,,\operatorname{window}_X,
                 \lambda_0)+D_{hi}(e,Y[0:o),\lambda_0).
\]
The inactive caches describe current words; the active cache describes its
frozen phase-entry word, with subsequent active writes represented physically.
An inactive matched tape is refreshed from its first current matching pair
at every positive time. Initially every cache is BLANK with zero error.
An inactive write $u\to v$ at $y$ changes its error by
$\mathbf1_{\{y\}}(u+v)$, leaving its base fixed.

We verify the support bound at the episode boundary. If an episode starts
at $u>0$, its tape was not active at $u-1$: an unmatched old carrier at
$u$ would have been retained. Maximality therefore says it was inactive
and matched at $u-1$. Its radius-$s$ refreshed error (or exact blank error
if $u-1=0$) acquires at most the one entering address. Subsequently its
support is contained in those old error addresses, that entering address,
and $D_h(I)$. Thus it is at most $s+1+m$, through every successful exit
write as well. An episode starting at zero needs only $D_h(I)$. A matched
inactive tape promoted immediately after one write needs at most $s+1$.
These cases cover every inactive cache before refresh or promotion.

Retain the lazy physical invariant, logical positions, current hashes and
continuous ports from \cref{thm:reference-carrier}. Reconstruct inactive
symbols directly from their caches, even when currently unmatched. The
active symbol comes from the physical tape on its visited interval and
from its frozen cache elsewhere. All queries use the current logical-head
raw ports and the same old windows/seed. Reconstruct the complete old tuple,
then compute the source rule, observer/hash updates and tentative error
updates for every old inactive tape. Keep the old active cache frozen.

After computing the optional append, test current matches and choose the
new carrier by the canonical policy. If it stays, extend the visited
interval. If it changes to $g$, freeze $g$'s exact tentative POST-write
cache, set the new offset, and reset the visited interval. The previous
carrier must now be matched, so refresh its inactive cache from its new
matching pair; refresh all other matched inactive tapes, and preserve the
remaining updated historical caches. Only the old carrier makes the one
physical write and move. The support argument proves that every cache
needed at a switch is legal, including a blank base with nonzero error.

INIT installs public original data, zero actual hashes and blank caches,
and obtains initial raw hashes from the one joint seed/hash label, within
the first transition evaluation. Multiple unmatched tapes at time zero
cause no problem because all actual words are blank. On invalid branches
finite range/support checks can halt. The invariant proves actual source
semantics through halt on the good inputs without testing the promise.
No episode counter, old-reference search, fresh label, clearing or seek is
needed. Each of the $k$ caches has
$[1+(T+1)(P+1)^2\ell]K_M$ descriptions. Multiply by the same state,
observer, head, interval and hash fields as before to obtain the claimed
product, including the active frozen slot. The label bill is unchanged.
\end{informalproof}

\begin{corollary}[Many distinct secondary modifications]\label{cor:secondary-modifications}
\partformal{ReferenceTail.lean}{cache_tail_of_half_simulation,reference_cache_sparse_reserves}{The exact budgets and success tail are checked from one common native-prefix half-set simulation with the stated cardinality bills. The historical-cache compiler and its canonical-episode invariant remain unformalized.}
Put $N=\lfloor640apC^2\rfloor+1$ and $A_0=k(\ell+2)$. Assume
\[
 q_{\rm src}+q_{\rm obs}\le P/(128N^2),\quad
 \lceil\log_2K_M\rceil\le P/(256A_0N^2),\quad
 Z_0\le P/(128N),\quad r_0+1\le\lfloor P/(8pN)\rfloor.
\]
Uniformly for every $b\ge1$ at sufficiently large $d$, fewer than
$2^{1+(1-1/(8N))P}$ inputs both transpose by $T$ and satisfy every episode
bound $|D_h(I)|\le m$. There is a fixed $\delta>0$ for which one may
take $s\le\delta P$, $m=\lfloor\delta P\rfloor$. Outside that exception,
some secondary episode has at least $\lfloor\delta P\rfloor+1$ distinct
actual modifications, more than $\delta P$ successful steps, and work-head
span at least $\lfloor\delta P\rfloor$.
\end{corollary}
\begin{informalproof}
Write $\lambda_M=\lceil\log_2K_M\rceil$ and
$L_0=\lceil\log_2(2k(T+1)\ell)\rceil$. Since $K_s\le K_M$,
the compiler's bit bill is at most
\[
 q_{\rm src}+q_{\rm obs}+A_0\lambda_M+
       k(\ell+1)L_0+\lceil\log_2B_{\rm cache}\rceil+2.
\]
The last three terms are $O(\log(P+1))$, so the supplied reserves imply
$q_*\le P/(64N^2)$ and $k\ell r\le P/(128N^2)$ uniformly at large $d$.
Together with $Z_0$ this gives the label bound of
\cref{lem:multiport-transition}. Apply it and retain the half-set factor.
For the quantitative choice it suffices to take $\delta\le1/6$ and
\[
 3\delta\log_2\frac{e\lceil4C+1\rceil(2^a-1)}{3\delta}
       \le\frac1{512A_0N^2}.
\]
Eventually $M\le3\delta P$ and the sparse binomial bound, with its ceiling,
fits the reserve. Such $\delta>0$ exists. Negation of the episode promise
gives the distinct-address count. One successful step changes at most one
address on this tape, and $j$ distinct integer addresses have span at least
$j-1$, proving the remaining claims.
The checked \lean{ExecutionPrefix.many_episode_changes_duration_span} counts
actual changed addresses only before halt. It proves the strict duration
bound and the span bound even when the declared episode contains the halting
configuration; that configuration contributes no successful write.
\end{informalproof}

Throughout the episode its tape and the current carrier are two different
unmatched tapes; when $k>2$ the other identity may vary. This is not an
independence or output-usefulness conclusion. The shared static scan seed
conditions from \cref{sec:references} carry over without a seed per episode.
For static reference families the cache reduces to a larger-radius case;
for public or prefix-only historical bases, a polynomial catalogue enlargement
already suffices. The added interface here permits raw-dependent historical
bases while current masks and cutoffs evolve.
